\centerline{\bf \Huge Конструктивы и инварианты}

\centerline{\bf \Huge В алгебре}

1. Найдите какие-нибудь четыре попарно различных натуральных числа $a$, $b$, $c$, $d$, для которых числа  $a^2 + 2cd + b^2$  и  $c^2 + 2ab + d^2$  являются полными квадратами.

2. Существуют ли такие натуральные числа $a$, $b$, $c$, $d$, что 

$a^3+b^3+c^3+d^3=800$?


3. Существуют ли две положительные несократимые дроби со знаменателями, не превосходящими 100, сумма которых равна 86/111?

\centerline{\bf \Huge В геометрии}

4. На прямоугольном столе лежат несколько картонных прямоугольников. Их стороны параллельны сторонам стола. Размеры прямоугольников могут различаться, они могут перекрываться, но никакие два прямоугольника не могут иметь 4 общих вершин. Может ли оказаться, что каждая точка, являющаяся вершиной прямоугольника, является вершиной ровно трех прямоугольников?

5. Можно ли обклеить куб 12 равными квадратами без разрезаний и наложений? (обклеить - покрыть все точки куба)



\centerline{\bf \Huge В числах}

6. Можно ли расставить в вершинах и на ребрах куба числа 1,2,$\ldots$, 20 так, чтобы на ребре стояло число, равное сумме чисел в вершинах этого ребра?

7. Существуют ли такие натуральные числа $a$, $b$ и $c$, что $a^3+b^3+c^3 = 500$?

8. Какое семизначное число является наименьшим, составленным из разных цифр, которое делится на 21?

9. Квадрат $2018 \times 2018$ разделен на квадратики размером $4\times4$, $3\times3$, $2\times2$ и $1\times1$. Может ли оказаться суммарное число квадратов $4\times4$, $2\times2$ и $1\times1$ равно 2018?

10. Можно ли расставить числа $1,2,3, \ldots, 26$ в вершинах и на сторонах тринаддцатиугольника так, чтобы сумма трех чисел на каждой стороне была одинаковой?

\centerline{\bf \Huge В шахматах}


11. Расставьте 12 ферзей на шахматной доске $8\times8$ так, чтобы каждый бил ровно трех других (ферзь - фигура, которая ходит по вертикали, горизонтали и диагонали).

12. Какое наибольшее количество ферзей можно расставить на шахматной доске $8\times8$ так, чтобы каждая из этих фигур была под ударом не более чем одной из остальных? 

13. Саша выставляет на пустую шахматную доску ладьи: первую – куда захочет, а каждую следующую ставит так, чтобы она побила нечётное число ранее выставленных ладей. Какое наибольшее число ладей он сможет так выставить?


\textit{Чудо формула с вебинара.} Докажите, что для любого $n$ натурального выполнено равенство:

$1^3 + 2^3 + \dotsc + (n-1)^3 + n^3 = (1 + 2 + \dotsc + n-1 + n)^2$

\textit{Решение: Первый способ.} \textbf{Индукция по $n$:}

\textit{База индукции ($n = 1$):}

При $n = 1$, левая и правая части равенства равны:

Левая часть: $1^3 = 1$

Правая часть: $(1)^2 = 1$

Таким образом, база индукции верна.

\textit{Индукционный переход:}

Предположим, что утверждение верно для $n = k$, т.е.:

\[
1^3 + 2^3 + \dots + (k-1)^3 + k^3 = (1 + 2 + \dots + k-1 + k)^2
\]

Теперь докажем, что утверждение также верно для $n = k + 1$:

\[
1^3 + 2^3 + \dots + (k-1)^3 + k^3 + (k + 1)^3
\]

Мы знаем, что

\[
1^3 + 2^3 + \dots + (k-1)^3 + k^3 = (1 + 2 + \dots + k-1 + k)^2
\]

(это следует из предположения индукции).

Теперь добавим $(k + 1)^3$ к обеим сторонам:

\[
(1 + 2 + \dots + k-1 + k)^2 + (k + 1)^3
\]

Это можно переписать как $((1 + 2 + \dots + k-1 + k) + (k + 1))^2$.

Упростим выражение внутри квадрата:
\[
((1 + 2 + \dots + k-1 + k) + (k + 1))^2 = (1 + 2 + \dots + k-1 + k + k + 1)^2
\]

Таким образом, утверждение верно и для $n = k + 1$.

Так как утверждение верно для $n = 1$ и если оно верно для $n = k$, то оно также верно для $n = k + 1$, по принципу математической индукции мы можем сказать, что данное равенство выполняется для любого $n$ натурального числа.
 \begin{flushright}
     Ч.Т.Д.
 \end{flushright}

 \textit{Второй способ.} А вторым способом будет который придумает кто-то из вас. Лучшее решение будет поощрено.

